% GENERATED FILE. Edit canonical lessons, then run npm run generate:books.
% canonical-source-hash: fnv1a64:bf837ad5ea72cc1c
% canonical-lessons: ES-C684-forastero, ES-C684-urbano, ES-C684-rural, ES-C684-residir, ES-C684-vecindario

\chapter{Ciudad o campo}
\label{ch:es-ciudad-o-campo}
\hlchaptermodality{684}

\begin{chapteropening}
\textbf{By the end of this chapter:} \emph{I can compare city and country life and talk about locals and people from out of town.}

Complete the last lesson of chapter 684: I can compare city and country life and talk about locals and people from out of town.
\end{chapteropening}

\section[el forastero]{el forastero — the outsider, the stranger (from out of town)}

\label{lesson:ES-C684-forastero}



\begin{quote}
Before the new one: say the Spanish for the hamlet, then the Spanish for the mayor’s office.
\end{quote}

\subsection*{You'll want to know: el forastero}
\textbf{el forastero} — \textquotedblleft{}the outsider, the stranger (from out of town)\textquotedblright{}.

\begin{quote}
En el pueblo todos saben quién es forastero.
\end{quote}

\begin{quote}
\emph{In the village everyone knows who is from out of town.}
\end{quote}

As an adjective: \emph{gente forastera}. It is not quite \emph{extranjero}, a foreigner.

\begin{cousinweb}
\emph{forastero} was borrowed from Catalan \emph{foraster} (EtymDB), built on \emph{fora} \textquotedblleft{}outside\textquotedblright{}, from Latin \emph{forās} \textquotedblleft{}out of doors\textquotedblright{}.
\end{cousinweb}

\begin{grammarlens}[title={What you've built}]
One more word for telling, explaining and describing at B1.
\end{grammarlens}

\subsection*{Your turn}
\begin{itemize}
\raggedright
  \item \emph{Say it:} \emph{el forastero}
  \item \emph{Say it:} \emph{el forastero}, once more
  \item \emph{Say it:} a sentence of your own with \emph{el forastero}
  \item \emph{Recall:} say the Spanish for the hamlet, then the Spanish for the mayor’s office, then say \emph{el forastero} again
\end{itemize}

\begin{tcolorbox}[breakable,colback=teal!4,colframe=teal!35!black,title={Before you move on}]
What does el forastero mean? (\textquotedblleft{}The outsider, the stranger (from out of town)\textquotedblright{}.) Say it once more.
\end{tcolorbox}
\section[urbano]{urbano — urban, city (as an adjective)}

\label{lesson:ES-C684-urbano}



\begin{quote}
Before the new one: say the Spanish for the mayor’s office, then the Spanish for the outsider.
\end{quote}

\subsection*{You'll want to know: urbano}
\textbf{urbano} — \textquotedblleft{}urban, city (as an adjective)\textquotedblright{}.

\begin{quote}
El autobús urbano es muy barato.
\end{quote}

\begin{quote}
\emph{The city bus is very cheap.}
\end{quote}

\emph{Transporte urbano} — city transport.

\begin{cousinweb}
\emph{urbano} is a learned borrowing from Latin \emph{urbānus} \textquotedblleft{}of the city\textquotedblright{} (EtymDB), from \emph{urbs} \textquotedblleft{}city\textquotedblright{}.
\end{cousinweb}

\begin{grammarlens}[title={What you've built}]
One more word for telling, explaining and describing at B1.
\end{grammarlens}

\subsection*{Your turn}
\begin{itemize}
\raggedright
  \item \emph{Say it:} \emph{urbano}
  \item \emph{Say it:} \emph{urbano}, once more
  \item \emph{Say it:} a sentence of your own with \emph{urbano}
  \item \emph{Recall:} say the Spanish for the mayor’s office, then the Spanish for the outsider, then say \emph{urbano} again
\end{itemize}

\begin{tcolorbox}[breakable,colback=teal!4,colframe=teal!35!black,title={Before you move on}]
What does urbano mean? (\textquotedblleft{}Urban, city (as an adjective)\textquotedblright{}.) Say it once more.
\end{tcolorbox}
\section[rural]{rural — rural, country (as an adjective)}

\label{lesson:ES-C684-rural}



\begin{quote}
Before the new one: say the Spanish for the outsider, then the Spanish for urban.
\end{quote}

\subsection*{You'll want to know: rural}
\textbf{rural} — \textquotedblleft{}rural, country (as an adjective)\textquotedblright{}.

\begin{quote}
Pasamos una semana en una casa rural.
\end{quote}

\begin{quote}
\emph{We spent a week in a country cottage.}
\end{quote}

In Spain \emph{una casa rural} is a holiday cottage in the country.

\begin{cousinweb}
\emph{rural} is a learned borrowing from Latin \emph{rūrālis} \textquotedblleft{}of the countryside\textquotedblright{} (EtymDB), from \emph{rūs} \textquotedblleft{}the country\textquotedblright{}.
\end{cousinweb}

\begin{grammarlens}[title={What you've built}]
One more word for telling, explaining and describing at B1.
\end{grammarlens}

\subsection*{Your turn}
\begin{itemize}
\raggedright
  \item \emph{Say it:} \emph{rural}
  \item \emph{Say it:} \emph{rural}, once more
  \item \emph{Say it:} a sentence of your own with \emph{rural}
  \item \emph{Recall:} say the Spanish for the outsider, then the Spanish for urban, then say \emph{rural} again
\end{itemize}

\begin{tcolorbox}[breakable,colback=teal!4,colframe=teal!35!black,title={Before you move on}]
What does rural mean? (\textquotedblleft{}Rural, country (as an adjective)\textquotedblright{}.) Say it once more.
\end{tcolorbox}
\section[residir]{residir — to reside, to live}

\label{lesson:ES-C684-residir}



\begin{quote}
Before the new one: say the Spanish for urban, then the Spanish for rural.
\end{quote}

\subsection*{You'll want to know: residir}
\textbf{residir} — \textquotedblleft{}to reside, to live\textquotedblright{}.

\begin{quote}
Residimos en la ciudad desde hace años.
\end{quote}

\begin{quote}
\emph{We have lived in the city for years.}
\end{quote}

A formal word for \emph{vivir}, common on forms.

\begin{cousinweb}
\emph{residir} is a learned borrowing from Latin \emph{residēre} \textquotedblleft{}to remain, to settle\textquotedblright{} (EtymDB), built on \emph{sedēre} \textquotedblleft{}to sit\textquotedblright{}.
\end{cousinweb}

\begin{grammarlens}[title={What you've built}]
One more word for telling, explaining and describing at B1.
\end{grammarlens}

\subsection*{Your turn}
\begin{itemize}
\raggedright
  \item \emph{Say it:} \emph{residir}
  \item \emph{Say it:} \emph{residir}, once more
  \item \emph{Say it:} a sentence of your own with \emph{residir}
  \item \emph{Recall:} say the Spanish for urban, then the Spanish for rural, then say \emph{residir} again
\end{itemize}

\begin{tcolorbox}[breakable,colback=teal!4,colframe=teal!35!black,title={Before you move on}]
What does residir mean? (\textquotedblleft{}To reside, to live\textquotedblright{}.) Say it once more.
\end{tcolorbox}
\section[el vecindario]{el vecindario — the neighbourhood; the residents}

\label{lesson:ES-C684-vecindario}



\begin{quote}
Before the new one: say the Spanish for rural, then the Spanish for to reside.
\end{quote}

\subsection*{You'll want to know: el vecindario}
\textbf{el vecindario} — \textquotedblleft{}the neighbourhood; the residents\textquotedblright{}.

\begin{quote}
El vecindario es muy tranquilo.
\end{quote}

\begin{quote}
\emph{The neighbourhood is very quiet.}
\end{quote}

Common in Latin America; Spain often says \emph{el barrio}.

\begin{cousinweb}
\emph{vecindario} is built on \emph{vecino} \textquotedblleft{}neighbour\textquotedblright{}, inherited from Latin \emph{vīcīnus} (EtymDB), from \emph{vīcus} \textquotedblleft{}a village, a street\textquotedblright{}.
\end{cousinweb}

\begin{grammarlens}[title={What you've built}]
One more word for telling, explaining and describing at B1.
\end{grammarlens}

\subsection*{Your turn}
\begin{itemize}
\raggedright
  \item \emph{Say it:} \emph{el vecindario}
  \item \emph{Say it:} \emph{el vecindario}, once more
  \item \emph{Say it:} a sentence of your own with \emph{el vecindario}
  \item \emph{Recall:} say the Spanish for rural, then the Spanish for to reside, then say \emph{el vecindario} again
\end{itemize}

\begin{tcolorbox}[breakable,colback=teal!4,colframe=teal!35!black,title={Before you move on}]
What does el vecindario mean? (\textquotedblleft{}The neighbourhood; the residents\textquotedblright{}.) Say it once more.
\end{tcolorbox}
